Los antecedentes de esta investigación provienen de una revisión sistematizada de catorce trabajos
(Fase~0, Grupo~9.2), agrupados en seis ejes (E1--E6) con fichas en el Anexo~B. La síntesis
extendida aparece en las secciones~\ref{sec:estado-arte} y~\ref{sec:brecha}; aquí se condensan las
cuatro líneas que convergen hacia el diseño del modelo de alineación.

\begin{enumerate}
  \item \textbf{Madurez de OpenAPI.} Casas et al.\ (2021), en un mapeo sistemático de 47 artículos
        (2011--2020), confirman que OpenAPI/Swagger es el estándar de facto para describir APIs REST
        y que su uso trasciende la documentación ---generación, pruebas, gobernanza---. Ninguno de
        esos trabajos aborda la coherencia con modelos de negocio bancarios de referencia como BIAN.
  \item \textbf{Correspondencia estructural y semántica.} Shvaiko y Euzenat (2005) consolidan el
        marco de \textit{schema matching}; El Bayadh et al.\ (2015), Chandrasekaran y Mago (2022) y
        Santos Sousa et al.\ (2026) cubren alineación ontológica y similitud semántica. Son técnicas
        aplicables por separado a componentes del contrato OpenAPI y del \ServiceDomain{} BIAN, pero
        ninguna las integra en un score único sobre el par bancario.
  \item \textbf{Separación dominio--interfaz.} Schmidt y Thiry (2020), en una SLR de 27 estudios
        sobre identificación de microservicios con MDE y DDD, muestran cómo las capacidades de
        dominio se traducen en límites de servicio. Ese marco justifica contrastar modelos de negocio
        e interfaces técnicas sin confundir la capa REST con el \ServiceDomain{} BIAN completo.
  \item \textbf{Delimitación del foco.} Las revisiones de testing, seguridad y selección de servicios
        (ejes E5 y E6) emplean OpenAPI con fines operativos ---cobertura de código, vulnerabilidades,
        ranking por QoS--- distintos a la alineación documental; acotan el alcance del presente plan
        (véase \Cref{sec:alcances}).
\end{enumerate}

De estas cuatro líneas se desprende que existen bases parciales para medir correspondencia
estructural (\StructuralScore{}), similitud semántica (\SemanticScore{}) y cobertura conceptual
(\CoverageScore{}), pero \textbf{no hay un procedimiento integrado} que los combine en un
\AlignmentScore{} reproducible sobre contratos OpenAPI bancarios frente a \ServiceDomain{}s BIAN. La
revisión detallada por ejes, las tablas sistematizadas y la formulación explícita de la brecha se
desarrollan a continuación en este capítulo.
